%!TEX root = ../main.tex

\chapter{Introducción}
\label{ch:introduccion}

\section{Motivación}
\label{sec:motivacion}

En el ámbito de los vehículos aéreos no tripulados de relativo bajo coste, como
drones de uso personal, los controladores como Betaflight o iNav ya han resuelto
el problema de la fusión de sensores en hardware de bajo coste. En estos sistemas
se usan sensores y microcontroladores muy populares y algoritmos relativamente
fáciles de implementar en software. El hardware y software se encuentran en un
estado muy maduro. El problema es que para uso educativo estos sistemas utilizan software con millones de líneas de código que no los hacen intuitivos para
estudiantes. Además, el PFD en estos sistemas no sigue las regulaciones presentes
en normativa aeronáutica, lo que tampoco los hace deseables para uso didáctico.

Por otro lado, en las clases de ingeniería los estudiantes estudian
la instrumentación de cabina de aeronaves (indicadores de actitud, altímetros,
etc.) a nivel teórico o mediante simuladores. Estos simuladores permiten
familiarizarse con la interfaz de un PFD, pero detrás no hay hardware real, lo
que los hace menos intuitivos y más difícil de aprender de ellos.

De estos dos factores nace la idea de este proyecto: una plataforma que conecte
ambos mundos, usando el mismo tipo de hardware y algoritmos del ecosistema de
drones para alimentar un PFD basado en las directrices de la regulación aeronáutica. De
esta forma, el estudiante puede recorrer el sistema completo de la
instrumentación aviónica, desde la lectura del sensor hasta la representación en
pantalla, incluyendo el hardware físico y el software. Además, el sistema
completo se puede replicar por menos de \SI{100}{\eur} en materiales, haciéndolo
fácil de incorporar en laboratorios, y todo el software empleado es de código
abierto, lo que permite que otros estudiantes o docentes puedan reproducirlo sin
coste de licencias.


\section{Objetivo}
\label{sec:objetivo}

El objetivo de este proyecto es el diseño e implementación de una plataforma
experimental de bajo coste orientada a la emulación de instrumentación
aviónica. El sistema se basa en un microcontrolador como núcleo y un
conjunto de sensores que incluyen sensores inerciales, barómetro, GPS, radar y
telemetría inalámbrica. A partir de la información proporcionada por los
sensores se calcula la actitud, el rumbo, la altitud y la velocidad vertical de
la plataforma en tiempo real. Sobre estos datos se desarrolla el software de
visualización en una pantalla similar a la de un \textit{Primary Flight Display}
(PFD).

El resultado final es un sistema integrado por tres subsistemas principales:

\begin{enumerate}
    \item Una PCB que integra los componentes físicos del proyecto, como el
    controlador y los sensores.
    \item El firmware que recibe los datos de los sensores, los integra mediante
    fusión sensorial y transmite los resultados al tercer subsistema.
    \item Un PFD implementado en el lenguaje Processing que muestra los valores
    calculados.
\end{enumerate}


\section{Alcance del proyecto}
\label{sec:alcance}

El alcance del proyecto se ha organizado en los siguientes paquetes de trabajo,
que estructuran también el contenido de la memoria:

\begin{enumerate}
    \item \textbf{Estado del arte:} estudio del estado actual en
    instrumentación aviónica, sistemas AHRS y visualización de vuelo, junto con
    el análisis de las normativas aplicables.

    \item \textbf{Definición de requisitos y arquitectura del sistema:} a partir
    del estado del arte se establecen los requisitos funcionales, de rendimiento,
    restricciones y normativos del sistema, y se define la arquitectura general en
    bloques funcionales, que orientan el resto del diseño.

    \item \textbf{Selección de componentes:} selección de los componentes del
    sistema físico, principalmente el controlador y los sensores, basada en el
    estudio del estado del arte.

    \item \textbf{Diseño PCB:} diseño del esquema eléctrico y la placa de
    circuito impreso que integra todos los módulos del sistema, microcontrolador
    y sensores seleccionados.

    \item \textbf{Firmware y calibración:} desarrollo de la arquitectura del
    firmware en Arduino, implementación de la calibración de los sensores y
    lectura de los datos.

    \item \textbf{Filtro de fusión de sensores (AHRS):} selección del filtro más
    adecuado en función de varios parámetros e implementación del mismo para
    obtener los parámetros deseados a partir de los valores de los sensores.

    \item \textbf{Sistema de telemetría:} desarrollo del sistema de transmisión
    inalámbrica entre el transmisor (microcontrolador y sensores) y el receptor
    (PC).

    \item \textbf{Visualización PFD:} implementación del Primary Flight Display
    completo con la disposición estándar Basic-T, marcada por la regulación
    correspondiente, y desarrollada en el entorno Processing.

    \item \textbf{Integración del sistema y pruebas de validación:} integración
    de todos los subsistemas desarrollados (hardware, firmware, telemetría y
    PFD) y documentación de los resultados obtenidos.
\end{enumerate}


\section{Actividades excluidas y limitaciones}
\label{sec:exclusiones}

Las siguientes actividades quedan fuera del alcance de este proyecto, aunque se
pueden identificar como líneas de trabajo futuro:

\begin{itemize}
    \item \textbf{Integración en una plataforma de vuelo.} La plataforma está
    diseñada para operar en tierra de forma experimental, no a bordo de un dron
    o aeronave real.

    \item \textbf{Representación del \textit{Flight Path Vector}.} Aunque sería una
    adición interesante para el PFD, requeriría datos de viento y navegación
    que están fuera del alcance del proyecto.

    \item \textbf{Certificación aeronáutica.} Aunque se sigue una serie de regulaciones
    para el desarrollo de algunos subsistemas, el resultado final está pensado
    para uso educativo y no aspira a ningún tipo de certificación.

    \item \textbf{Sistemas aviónicos adicionales.} No se contempla la incorporación
    de sistemas como radar meteorológico, TCAS, ADS-B o FMS.
\end{itemize}

Como limitación principal, el sistema no está diseñado ni para certificación
aeronáutica ni para el control real de un vehículo aéreo. Por este motivo no
incorpora redundancia, protección ambiental ni los sistemas de seguridad que
serían requeridos para un sistema real. Esto es coherente con su propósito: se
trata de una plataforma de aprendizaje, no de un equipo de vuelo.